% ==============================================================
\section{はじめに}\label{sec:intro}
% ==============================================================

頻出パターンマイニング~\cite{agrawal1994fast,han2000mining,fournier2017survey}は，
トランザクションデータ中に繰り返し出現するアイテムセットを発見する基盤的技術であり，
小売分析，Webマイニング，バイオインフォマティクスなど幅広い分野で応用されている．
Apriori~\cite{agrawal1994fast}，FP-Growth~\cite{han2000mining}，ECLAT~\cite{zaki2000scalable}
といった古典的手法はグローバル頻出パターンを効率的に列挙し，
スライディングウィンドウ型手法~\cite{dsaa2025}は
特定時間区間内で局所的に密集するパターンを検出する．

しかし，これらの手法は\emph{どの}パターンが\emph{いつ}頻出であるかを同定するのみであり，
\emph{なぜ}パターンの出現頻度が変動するかを説明する機能を持たない．
実際の応用場面では，販促キャンペーン，祝日，規制変更，システム更新といった
外部イベントがパターンサポートの時間的変動を駆動していることが多い．
「どのイベントがどのパターンの変動を引き起こしたか」を特定できれば，
キャンペーン効果の定量化，副作用の検出，システム異常の診断など，
意思決定に直結する知見が得られる．

この帰属タスクは，以下の3点で技術的に困難である．
第一に，\textbf{仮説の爆発}:
$|\mathcal{F}|$ 個のパターンと $|\mathcal{E}|$ 個のイベントの全組合せ
$O(|\mathcal{F}| \times |\mathcal{E}|)$ を検定する必要があり，
多重検定補正により検出力が急速に低下する．
第二に，\textbf{副次パターンの連鎖}:
イベントがアイテム $i$ の頻度を増加させると，
$i$ を含む全パターンのサポートが相関して変動し，
統計検定だけでは除去できない大量の偽陽性が発生する．
第三に，\textbf{サポート時系列の分布自由性}:
パターンサポートの時系列分布はパターンごとに異なり，
パラメトリックな仮定を置くことが困難である．

このように帰属タスクは実用上重要であるにもかかわらず，
頻出パターンのサポート変動を外部イベントに結びつける統計フレームワークは，
我々の知る限り存在しない．
時間的パターンマイニング~\cite{mahanta2005finding,tanbeer2009discovering,kiran2017discovering,fournier2021mining}
は周期性や局所頻度を持つパターンを発見するが，
パターン変動と外生イベントの関連はモデル化しない．
変化点検出手法~\cite{page1954continuous,aminikhanghahi2017survey}は
個別時系列の構造変化を同定するが，
数千のパターンサポート時系列を同時に検定する多重検定の負荷や，
検出された変化をイベントに紐づける機能は備えていない．
計量経済学のイベント・スタディ法~\cite{mackinlay1997event}は
企業イベントの株価への影響を評価するが，
単一のアウトカム変数を対象としており，
数千のパターン時系列を同時に扱う設定には適用できない．
ベイズ構造時系列モデルに基づく CausalImpact~\cite{brodersen2015causalimpact}や
イベント一致分析~\cite{donges2016event}は個別の時系列に対しては有効だが，
反単調性に起因する大量の相関パターンを同時に扱う設計にはなっていない．

\subsection{本研究の貢献}

本稿では，頻出パターンマイニングと外部イベント帰属を橋渡しする
統計フレームワーク\textbf{「イベント帰属パイプライン」}を提案する．
貢献は以下の3点である．

\begin{enumerate}
\item \textbf{問題定式化と5ステップパイプライン}（\ref{sec:problem}--\ref{sec:method}節）:
パターンサポート時系列と外部イベントの帰属問題を形式的に定義し，
密集区間からの変化点抽出，変化量の評価，帰属スコアリング，
円形シフト置換検定によるグローバル Benjamini--Hochberg~\cite{benjamini1995controlling} FDR補正，
およびアイテム重複排除からなる統合フレームワークを提案する．
個々の構成要素は既存技術に基づくが，
パターンマイニング固有の課題（副次パターン問題，仮説爆発）に対処する
二段階スクリーニング設計（統計検定＋構造的重複排除）が本手法の核心である．

\item \textbf{副次パターン問題の同定と解決}（\ref{sec:dedup}節）:
真に影響を受けたパターンとアイテムを共有する副次パターンが
相関した偽陽性を生む問題（\emph{副次パターン問題}）を定義し，
これが統計検定単独では解消不可能であることを理論的に示す．
イベントごとの Union-Find クラスタリングに基づく重複排除で解決する．

\item \textbf{包括的な実験評価}（\ref{sec:experiments}節）:
(i) 合成データにおける完全検出（F1 = 1.0），
(ii) パラメータ選択に対する頑健性，
(iii) 100万トランザクションまでの線形スケーラビリティ，
(iv) 実データにおける偽陽性96\%削減を実証する．
同時に，高ベースラインサポートパターンでの検出失敗や
イベント密集時の精度低下など，手法の限界も明確に報告する．
\end{enumerate}
